\section{哈代域的扩张}

给定一个哈代域 \(\mathfrak{K}\) ，我们将看到如何构造新的哈代域 \(\mathfrak{K}' \supset \mathfrak{K}\) ，使 \(\mathfrak{K}' / \mathrm{R}_{\infty}\) 是通过（在代数意义上，cf.~Alg., V, §2）向 \(\mathfrak{K} / \mathrm{R}_{\infty}\) 添加我们将明确其形状的新元素而得到的。

引理 1. 设 \(a(x), b(x)\) 是在区间 \([x_0, +\infty[\) 上不变号的两个连续实函数。若函数 \(y(x)\) 在该区间上连续、可导且满足恒等式

\[
y ^ {\prime} = a y + b\tag{1}
\]

则存在一个区间 \([x_{1}, +\infty[\) ，y 在其上不变号。

事实上，令 \(z(x) = y(x) \exp \left(-\int_{x_0}^x a(t)dt\right)\) (cf.~IV, p.~183)；则由 (1)， \(z'(x) = b(x) \exp \left(-\int_{x_0}^x a(t)dt\right)\) 。若 \(b(x) \geqslant 0\) （对 \(x \geqslant x_0\) ），则 \(z\) 在该区间上递增，于是或者在整个区间上 \(< 0\) ，或者在某个区间 \([x_1, +\infty[\) 上为零，或者 \(> 0\) （在某个区间 \([x_1, +\infty[\) 上）；由于 \(y\) 与 \(z\) 同号，命题在这一情形得证。若 \(b(x) \leqslant 0\) （对 \(x \geqslant x_0\) ），论证类似。

注. 这个非常初等的性质不能推广到阶 \textgreater{} 1 的线性微分方程；例如，函数 \(y = \sin x\) 满足 \(y'' + y = 0\) ，但在 \(+\infty\) 的每个邻域上都变号。

引理 2. 设 \(a(x)\) 与 \(b(x)\) 是属于给定哈代域 \(\mathfrak{K}\) 的两个函数， \(y(x)\) 是在区间 \([x_0, +\infty[\) （ \(a\) 与 \(b\) 在其上有定义且连续）上满足恒等式 (1) 的函数。若 \(p(u)\) 是 \(u\) 的一个多项式，其系数是 \(x\) 的、属于 \(\mathfrak{K}\) 的函数，在 \([x_0, +\infty[\) 上有定义且可导，则存在一个区间 \([x_1, +\infty[\) ，函数 \(p(y)\) 在其上不变号。

若 \(p(u)\) 的系数在 \([x_{0}, +\infty[\) 上恒为零，或者 \(p(u)\) 关于 u 是 0 次的，则命题是平凡的，因为 K 中任何函数都在某个区间 \([x_{1}, +\infty[\) 上具有常号。设 \(p(u)\) 是 n \textgreater{} 0 次的；于是 \(p(u)\) 的首项系数 c 为 \(\neq 0\) （在一个区间 \([\alpha, +\infty[\) 上）；从而可写 \(p(u) = c(u^{n} + c_{1}u^{n-1} + \cdots + c_{n})\) ，其中 \(c, c_{1}, c_{2}, \ldots, c_{n}\) 是属于 K 且在 \([\alpha, +\infty[\) 上可导的函数；故只需对 c = 1 证明引理。我们对 n 作归纳论证；有

\[
\begin{array}{c} \frac {d}{d x} (p (y)) = (a y + b) \bigl (n y ^ {n - 1} + (n - 1) c _ {1} y ^ {n - 2} + \dots + c _ {n - 1} \bigr) \\ + c _ {1} ^ {\prime} y ^ {n - 1} + \dots + c _ {n} ^ {\prime} = n a p (y) + q (y) \end{array}
\]

其中 \(q(u)\) 是次数 \(\leqslant n - 1\) 、系数在 \(\mathfrak{K}\) 中的多项式。按假设，函数 \(na(x)\) 与 \(q(y(x))\) 在某个区间 \([\beta, +\infty[\) 上不变号；于是本引理是引理 1 的推论。

定理 1. 设 \(a(x)\) 与 \(b(x)\) 是属于给定哈代域 \(\mathfrak{K}\) 的两个函数， \(y(x)\) 是在区间 \([x_0, +\infty[\) 上满足 (1) 的函数。当 \(r(u) = p(u) / q(u)\) 取遍 \(u\) 的、系数在 \(\mathfrak{K}\) 中且使 \(q(y)\) 在 \(+\infty\) 的某个邻域上不恒为零的有理分式所成的集时，集 \(\mathfrak{K}(y)\) （由函数 \(r(y)\) 组成）构成一个哈代域。

事实上，由引理 2，存在一个区间 \([x_1, +\infty[\) ， \(r(y)\) 在其上有定义、连续且不变号，由此立即推出 \(\mathfrak{K}(y)/\mathrm{R}_{\infty}\) 是一个域；此外，由于

\[
\frac {d}{d x} (r (y)) = r ^ {\prime} (y) y ^ {\prime} = r ^ {\prime} (y) (a y + b)
\]

（其中 \(r'(y) = \left(p'(y)q(y) - p(y)q'(y)\right)/\left(q(y)\right)^2\) 按假设在 \(+\infty\) 的某个邻域上有定义）， \(\mathfrak{K}(y)\) 中每个函数的导数都属于 \(\mathfrak{K}(y)\) ，这就证明了 \(\mathfrak{K}(y)\) 满足 V, p.~247 的定义 1 的条件。

显然， \(\mathfrak{K}(y) / \mathbb{R}_{\infty}\) 是由 \(\mathfrak{K} / \mathbb{R}_{\infty}\) 通过 \(y\) 模 \(\mathbb{R}_{\infty}\) 的类的代数添加而得到的。也说 \(\mathfrak{K}(y)\) 是把 \(y\) 添加到 \(\mathfrak{K}\) 而得到的。

推论 1. 若 \(y\) 是 \(\mathfrak{K}\) 中的一个函数，在 \(+\infty\) 的某个邻域上不恒为零，则 \(\mathfrak{K}(\log |y|)\) 是一个哈代域。

事实上， \((\log |y|)' = y'/y\) 等于 \(\mathcal{K}\) 中的一个函数（在一个区间 \([x_0, +\infty[\) 上）。

推论 2. 若 \(y\) 是 \(\mathfrak{K}\) 中任一函数，则 \(\mathfrak{K}(e^y)\) 是一个哈代域。

事实上， \(\left(e^{y}\right)^{\prime} = e^{y}y^{\prime}\) ，且 \(y^\prime\) 等于 \(\mathcal{K}\) 中的一个函数（在一个区间 \([x_0, +\infty[\) 上）。

推论 3. 若 \(\mathfrak{K}\) 包含实常数，且 \(y\) 是 \(\mathfrak{K}\) 中的一个函数，在 \(+\infty\) 的某个邻域上不恒为零，则 \(\mathfrak{K}(|y|^{\alpha})\) 对每个实数 \(\alpha\) 都是一个哈代域。

事实上， \(\frac{d}{dx}\big(|y|^{\alpha}\big) = |y|^{\alpha}(\alpha y'/y)\) ，且 \(\alpha y'/y\) 等于 \(\mathfrak{K}\) 中的一个函数（在一个区间 \([x_0, +\infty[\) 上）。

最后我们注意，若 \(y\) 是 \(\mathfrak{K}\) 中任一函数的原函数，则 \(\mathfrak{K}(y)\) 仍是一个哈代域。

